\section{Discussion}
\label{sec:discussion}

\paragraph{One scalar, three operations.}
The practical appeal of \cref{eq:value} is consolidation: a single
learned value replaces the separate, independently tuned rules that
production systems use for write-time importance, eviction, and
retrieval ranking. A deployment exposes one seven-dimensional knob, and
\cref{eq:forget} and the encoding tiers inherit any change to it
automatically. Because the value is linear, the learned weights are a
\emph{readable audit} of the policy: \cref{sec:exp:blind} shows the
fitted policy says, in plain terms, ``keep reliable user-stated facts;
do not trust session-topic similarity for forgetting.''

\paragraph{The oracle/blind lesson generalises.}
Our sharpest methodological point is independent of the value function.
Any retention metric whose gold set is defined relative to a future
query will reward a policy that scores relevance \emph{to that query}
--- but that policy is unavailable at forgetting time. Reporting such
numbers measures an upper bound on retrieval, not the quality of
forgetting. Evaluations of agent forgetting should therefore (i) compute
relevance only from information available at consolidation, and (ii)
report the oracle ceiling alongside, so the retrieval/retention gap is
visible. The $0.979 \rightarrow 0.286$ collapse of goal-only between our
two regimes (\cref{tab:blind}) is the size of the artefact a careless
metric would hide.

\paragraph{Learned weights as a workload diagnostic.}
Because the weights are fit to the chosen objective, they characterise the
workload, not just the policy. On LongMemEval the signal is reliability
and self-relevance; a coding agent's workload might instead reward task
utility and goal relevance. The same harness, pointed at a new return
signal, yields a new interpretable weighting --- a portable recipe for
\emph{discovering} what a given agent should remember, rather than
guessing it. This extends the hand-set recency/importance/relevance blend
of Generative Agents \citep{park2023generative} in two ways: the blend is
learned, and it governs encoding and forgetting, not only retrieval.
